% ---
% id: receiver/ass:passive
% kind: ass
% title: "Passive linear optics"
% summary: >-
%   Every optical element is passive linear optics, with no gain and no mixing
%   of creation and annihilation operators. Total photon number is then a
%   conserved grading, on which every block statement of the module rests.
% prominence: standard
% status: accepted
% uses: []
% assumes: []
% refs: []
% provenance: newly-derived
% short_id: A-PASSIVE
% verification: never
% category: receiver
% affects: [secrecy]
% refines: A-TRUST
% user_statement: >-
%   Every optical component of your receiver is passive and linear: beam
%   splitters, wave plates, polarisers, attenuators and the like. No component
%   amplifies light, and none is nonlinear (no frequency conversion, no
%   squeezing). The total photon number is then never increased inside the
%   receiver. If a component can add photons, the photon-number blocks the
%   analysis rests on do not exist.
% lean:
% ---

Under this assumption the total photon number is a conserved grading of the
receiver's modes.

\begin{assumption}\label{ass:passive}
(i) Every optical element is passive linear optics: it maps annihilation
operators to combinations of annihilation operators,
$b_i=\sum_j\Gmat_{ij}a_j$, with $\Gmat^\dagger\Gmat\preceq\id$;
(ii) no element mixes $a$ with $a^\dagger$, and none provides gain.
\end{assumption}
